\subsection{Lattices}

DEFINITION 1. Let V be a finite-dimensional vector space over the field K. A lattice of V with respect to A (or simply a lattice of V) is defined to be any sub-A-module M of V satisfying the following condition:

There exist two free sub-A-modules \(L_{1}\) , \(L_{2}\) of V such that \(L_{1} \subset M \subset L_{2}\) and \(\operatorname{rg}_{\mathbf{A}}(\mathbf{L}_{1}) = \operatorname{rg}_{\mathbf{K}}(\mathbf{V})\) .

\textbf{Examples}\par

\begin{enumerate}
\def\labelenumi{(\arabic{enumi})}
\item
  If we take \(\mathbf{V} = \mathbf{K}\) , the lattices of \(\mathbf{K}\) are just the fractional ideals \(\neq (0)\) of \(\mathbf{K}\) ( \(\S 1\) , no. 1, Definition 1).
\item
  If \(\operatorname{rg}_{\mathbf{K}}(\mathbf{V}) = n\) , every free sub-A-module L of V has a basis containing at most n elements, every subset of V which is free over A being free over K; for L to be a lattice of V, it is necessary and sufficient that L have a basis of n elements (in other words, that \(\operatorname{rg}_{\mathbf{A}}(\mathbf{L}) = n\) ).
\item
  If A is a principal ideal domain, every lattice M of V is a finitely generated A-module (since A is Noetherian) which is torsion-free and hence a free A-module (Algebra, Chapter VII, § 4, no. 3, Corollary 2 to Theorem 2).
\end{enumerate}

PROPOSITION 1. For a sub-A-module M of V to be a lattice of V, it is necessary and sufficient that KM = V and that M be contained in a finitely generated sub-A-module of V.

The conditions are obviously necessary, for a free sub-A-module of V with the same rank as V generates V. Conversely, if KM = V, M contains a basis \((a_{i})_{1 \leqslant i \leqslant n}\) of V over K and hence it contains the free sub-A-module \(L_{1}\) generated by the \(a_{i}\) ; on the other hand, if \(M \subset M_{1}\) , where \(M_{1}\) is a sub-A-module of V generated by a finite number of elements \(b_{j}\) and \((e_{i})_{1 \leqslant i \leqslant n}\) is a basis of V over K, there exists an element \(s \neq 0\) of A such that each of the \(b_{j}\) is a linear combination of the \(s^{-1}e_{i}\) with coefficients in A; if \(L_{2}\) is the free sub-A-modules of V generated by the \(s^{-1}e_{i}\) , then \(M \subset L_{2}\) .

COROLLARY. Suppose that A is Noetherian; for a sub-A-module M of V to be a lattice of V, it is necessary and sufficient that KM = V and M be finitely generated.

Remark (1). Recall that for every sub-A-module M of V the canonical mapping \(M \otimes_{A} K \to V\) is injective and has image KM (Algebra, Chapter II, § 7, no. 10, Proposition 26); to say that KM = V means therefore that this mapping is bijective.

PROPOSITION 2. Let M be a lattice of V and \(M_{1}\) a sub-A-module of V. If there exist two elements x, y of \(K^{*}\) such that \(xM \subset M_{1} \subset yM\) , \(M_{1}\) is a lattice of V; conversely, if \(M_{1}\) is a lattice of V, there exist two non-zero elements a, b of A such that \(aM \subset M_{1} \subset b^{-1}M\) .

If \(L_{1}\) , \(L_{2}\) are two free lattices of V such that \(L_{1} \subset M \subset L_{2}\) , the relations \(xM \subset M_{1} \subset yM\) imply \(xL_{1} \subset M_{1} \subset yL_{2}\) and \(xL_{1}\) and \(yL_{2}\) are free lattices; conversely, if \(M_{1}\) is a lattice and \((e_{i})_{1 \leqslant i \leqslant n}\) a basis of \(L_{2}\) over A, the relation \(KM_{1} = V\) implies the existence of \(x = a/s \in K^{*}\) (where a and s are non-zero elements of A) such that \(xe_{i} \in M_{1}\) for all i, whence \(xM \subset xL_{2} \subset M_{1}\) and a fortiori \(aM \subset M_{1}\) ; exchanging the roles of M and \(M_{1}\) it can be similarly shown that there exists \(b \neq 0\) in A such that \(bM_{1} \subset M\) .

PROPOSITION 3. (i) If \(\mathbf{M}_1\) and \(\mathbf{M}_2\) are lattices of \(\mathbf{V}\) , so are \(\mathbf{M}_1 \cap \mathbf{M}_2\) and \(\mathbf{M}_1 + \mathbf{M}_2\) .

\begin{enumerate}
\def\labelenumi{(\roman{enumi})}
\setcounter{enumi}{1}
\item
  If \(\mathbf{W}\) is a vector subspace of \(\mathbf{V}\) and \(\mathbf{M}\) is a lattice of \(\mathbf{V}\) , \(\mathbf{M} \cap \mathbf{W}\) is a lattice of \(\mathbf{W}\) .
\item
  Let \(\mathbf{V},\mathbf{V}_1,\ldots ,\mathbf{V}_k\) be vector spaces of finite rank over \(\mathbf{K}\) and let
\end{enumerate}

\[
f \colon \mathrm{V} _ {1} \times \dots \times \mathrm{V} _ {k} \rightarrow \mathrm{V}
\]

be a multilinear mapping whose image generates V. If \(M_{i}\) is a lattice of \(V_{i}\) for \(1 \leqslant i \leqslant k\) , the sub-A-module of V generated by \(f(\mathbf{M}_{1} \times \cdots \times \mathbf{M}_{k})\) is a lattice of V.

\begin{enumerate}
\def\labelenumi{(\roman{enumi})}
\setcounter{enumi}{3}
\item
  Let \(\mathbf{V}\) and \(\mathbf{W}\) be two vector spaces of finite rank over \(\mathbf{K}\) , \(\mathbf{M}\) a lattice of \(\mathbf{V}\) and \(\mathbf{N}\) a lattice of \(\mathbf{W}\) . The sub-A-module \(\mathbf{N}:\mathbf{M}\) of \(\operatorname{Hom}_{\mathbf{K}}(\mathbf{V},\mathbf{W})\) , consisting of the K-linear mappings \(f\) such that \(f(\mathbf{M})\subset \mathbf{N}\) , is a lattice of \(\operatorname{Hom}_{\mathbf{K}}(\mathbf{V},\mathbf{W})\) .
\item
  By virtue of Proposition 2, there exist non-zero \(a\) and \(b\) in \(\mathbf{A}\) such that \(a\mathbf{M}_1 \subset \mathbf{M}_2 \subset b^{-1}\mathbf{M}_1\) ; we conclude that \(\mathbf{M}_1 \cap \mathbf{M}_2\) and \(\mathbf{M}_1 + \mathbf{M}_2\) lie between \(a\mathbf{M}_1\) and \(b^{-1}\mathbf{M}_1\) and are therefore lattices by virtue of Proposition 2.
\item
  Let S be a complement of W in V, \(L_{w}\) a free lattice of W and \(L_{s}\) a free lattice of S, so that \(L = L_{w} \oplus L_{s}\) is a free lattice of V. Then there exist x, y in \(K^{*}\) such that \(xL \subset M \subset yL\) . We deduce that \(xL_{w} \subset M \cap W \subset yL_{w}\) , which shows that \(M \cap W\) is a lattice of W (Proposition 2).
\item
  As \(KM_{i} = V_{i}\) , clearly by linearity \(f(\mathbf{M}_{1} \times \cdots \times \mathbf{M}_{k})\) generates the vector K-space V; on the other hand, for all i, there exists a finitely generated sub-A-module \(N_{i}\) of \(V_{i}\) such that \(M_{i} \subset N_{i}\) ; the sub-A-module N of V generated by \(f(\mathbf{N}_{1} \times \cdots \times \mathbf{N}_{k})\) is finitely generated and contains M and hence M is a lattice of V (Proposition 1).
\item
  Let P (resp. Q) be a free lattice of V (resp. W) containing M (resp. contained in N); obviously N: M ⊃ Q: P. Now it is immediate that Q: P is isomorphic to Hom \(_{A}\) (P, Q), hence is a free A-module of rank (rg \(_{A}\) P)(rg \(_{A}\) Q) (Algebra, Chapter II, § 1, no. 6, Corollary 1 to Proposition 6) and therefore a lattice of Hom \(_{K}\) (V, W). Similarly, if P' (resp. Q') is a free lattice of V (resp. W) contained in M (resp. containing N), then Q': P' ⊃ N: M and Q': P' is a lattice of Hom \(_{K}\) (V, W); whence the conclusion.
\end{enumerate}

\textbf{Remarks}\par

\begin{enumerate}
\def\labelenumi{(\arabic{enumi})}
\setcounter{enumi}{1}
\item
  Proposition 3 (i) shows that the set R(V) of lattices of V is lattice-ordered with respect to inclusion; moreover, if M is a fixed lattice of V, the xM, where x runs through K*, form a subset of R(V) which is both coinitial and cofinal (Set Theory, Chapter III, § 1, no. 7).
\item
  In the notation of Proposition 3 (iv), the canonical mapping
\end{enumerate}

\[
\mathrm{N}: \mathrm{M} \rightarrow \operatorname{Hom} _ {\mathrm{A}} (\mathrm{M}, \mathrm{N}),
\]

which maps every K-linear mapping \(f \in N\) : M to the A-linear mapping from M to N which has the same graph as \(f|_{M}\) , is bijective; for every A-linear mapping \(g\) : \(M \to N\) can be imbedded in a K-linear mapping

\[
g \otimes 1: M \otimes_ {A} K \rightarrow N \otimes_ {A} K
\]

and it has been seen that \(M \otimes_{A} K\) and \(N \otimes_{A} K\) are respectively identified with V and W.

In particular, if we take W = K, N = A, \(\operatorname{Hom}_{K}(V, W)\) is just the dual vector K-space \(V^{*}\) of V and A: M is identified with the dual A-module \(M^{*}\) of M; we shall henceforth make this identification and we shall say that \(M^{*}\) is the dual lattice of M: it is therefore the set of \(x^{*} \in V^{*}\) such that \(\langle x, x^{*} \rangle \in A\) for all \(x \in M\) .

COROLLARY. Let U, V, W be three vector spaces of finite rank over K and \(f: \mathbf{U} \times \mathbf{V} \to \mathbf{W}\) a left non-degenerate K-bilinear mapping (Algebra, Chapter IX, § 1, no. 1, Definition 3). If M is a lattice of V and N a lattice of W, the set N: \(f\mathbf{M}\) of \(x \in \mathbf{U}\) such that \(f(x, y) \in \mathbf{N}\) for all \(y \in \mathbf{M}\) is a lattice of U.

Let \(s_f \colon \mathbf{U} \to \operatorname{Hom}_{\mathbf{K}}(\mathbf{V}, \mathbf{W})\) be the K-linear mapping left associated with \(f\) (Algebra, Chapter IX, loc. cit.) such that \(s_f(x)\) is the linear mapping \(y \mapsto f(x, y)\) ; recall that to say that \(f\) is left non-degenerate means that \(s_f\) is injective. By Proposition 3 (iv), \(\mathbf{N} \colon \mathbf{M}\) is a lattice of \(\operatorname{Hom}_{\mathbf{K}}(\mathbf{V}, \mathbf{W})\) ; as \(\mathbf{N} \colon {}_f \mathbf{M} = s_f^{-1}(\mathbf{N} \colon \mathbf{M})\) and \(s_f\) is injective, the corollary follows from Proposition 3 (ii).

\textbf{Examples}\par

\begin{enumerate}
\def\labelenumi{(\arabic{enumi})}
\setcounter{enumi}{3}
\tightlist
\item
  Let S be a (not necessarily associative) K-algebra of finite rank with a unit element; then the bilinear mapping \((x,y)\mapsto xy\) of S × S to S is (left and right) non-degenerate. If M and N are lattices of S with respect to A, so are M.N (Proposition 3 (iii)) and the set of \(x\in S\) such that \(xM\subset N\) (Corollary to Proposition 3). Note that there exists a sub-A-algebra of S containing the unit element of S which is a lattice of S; for consider a basis \((e_{i})_{1\leqslant i\leqslant n}\) of S such that \(e_{1}\) is the unit element of S and let \(e_{i}e_{j}=\sum_{k}c_{ijk}e_{k}\) be the multiplication table of S ( \(1\leqslant i\leqslant n,1\leqslant j\leqslant n\) ), so that \(c_{1jk}=\delta_{jk},c_{11k}=\delta_{1k}\) (Kronecker symbols). Let \(s\in A\) be a non-zero and such that \(c_{ijk}^{\prime}=s.c_{ijk}\in A\) for all triplets of indices \((i,j,k)\) ; if we write \(e_{i}^{\prime}=s^{-1}e_{i}\) for \(i\geqslant2\) , then
\end{enumerate}

\[
e _ {i} ^ {\prime} e _ {j} ^ {\prime} = s c _ {i j 1} ^ {\prime} e _ {1} + \sum_ {k \geqslant 2} c _ {i j k} ^ {\prime} e _ {k} ^ {\prime}
\]

for \(i \geqslant 2\) and \(j \geqslant 2\) ; the lattice of S with basis \(e_1\) and the \(e_i' (2 \leqslant i \leqslant n)\) is a sub-A-algebra of S with unit element \(e_1\) .

\begin{enumerate}
\def\labelenumi{(\arabic{enumi})}
\setcounter{enumi}{4}
\tightlist
\item
  Let V be a finite-dimensional vector space over K and f a non-degenerate bilinear form on V. If M is a lattice of V, it follows from the Corollary to Proposition 3 that the set \(M_{f}^{*}\) of \(x \in V\) such that \(f(x, y) \in A\) for all \(y \in M\) is also a lattice of V; if \(s_{f}: V \to V^{*}\) is the linear mapping left associated with f (which is bijective), \(s_{f}(M_{f}^{*})\) is just the dual lattice \(M^{*}\) of M.
\end{enumerate}

PROPOSITION 4. Let B be an integral domain, A a subring of B and K and L the respective fields of fractions of A and B. Let V be a finite-dimensional vector space over K.

\begin{enumerate}
\def\labelenumi{(\roman{enumi})}
\item
  For every lattice \(\mathbf{M}\) of \(\mathbf{V}\) with respect to \(\mathbf{A}\) , the image \(\mathbf{BM}\) of \(\mathbf{M}_{(\mathbf{B})} = \mathbf{M} \otimes_{\mathbf{A}} \mathbf{B}\) in \(\mathbf{V}_{(\mathbf{L})} = \mathbf{V} \otimes_{\mathbf{K}} \mathbf{L}\) is a lattice of \(\mathbf{V}_{(\mathbf{L})}\) with respect to \(\mathbf{B}\) .
\item
  Suppose further that B is a flat A-module. Then the canonical mapping \(\mathbf{M}_{(\mathbf{B})}\to\mathbf{BM}\) is bijective. If further B is faithfully flat, the mapping which maps every lattice M of V with respect to A to the lattice BM of \(\mathbf{V}_{(\mathbf{L})}\) with respect to B is injective.
\item
  As \(\mathbf{KM} = \mathbf{V}\) , clearly \(\mathbf{L}\) . (BM) = \(\mathbf{V}_{(L)}\) ; on the other hand \(\mathbf{M}\) is contained in a finitely generated sub-A-module \(\mathbf{M}_1\) of \(\mathbf{V}\) and hence \(\mathbf{BM}\) is contained in \(\mathbf{BM}_1\) which is a finitely generated B-module; whence assertion (i) (Proposition 1).
\item
  \(V_{(L)} = V \otimes_{K} L = V \otimes_{A} L\) (Chapter II, § 2, no. 7, Proposition 18) and, as L is a flat B-module, it is also a flat A-module (Chapter I, § 2, no. 7, Corollary 3 to Proposition 8). As B is a flat A-module, the canonical mapping \(M \otimes_{A} B \to V \otimes_{A} B\) is injective; on the other hand, since V is a free K-module and K a flat A-module, V is a flat A-module (Chapter I, § 2, no. 7, Corollary 3 to Proposition 8) and hence the canonical mapping \(V \otimes_{A} B \to V \otimes_{A} L\) is injective, which establishes the first assertion. To see also that the relation \(BM_{1} = BM_{2}\) implies \(M_{1} = M_{2}\) for two lattices \(M_{1}, M_{2}\) of V with respect to A when B is a faithfully flat A-module, note first that \(BM_{1} \cap BM_{2} = B(M_{1} \cap M_{2})\) (Chapter I, § 2, no. 6, Proposition 6); we may therefore confine our attention to the case where \(M_{1} \subset M_{2}\) and our assertion then follows from Chapter I, § 3, no. 1, Proposition 3 applied to the canonical injection \(M_{1} \to M_{2}\) .
\end{enumerate}

COROLLARY. Suppose that A is a discrete valuation ring. Let \(\hat{\mathbf{A}}\) be its completion and let \(\hat{\mathbf{K}}\) be the field of fractions of \(\hat{\mathbf{A}}\) (Chapter VI, § 5, no. 3). The mapping \(\phi\) , which maps every lattice M of V to the lattice \(\hat{\mathbf{A}}\mathbf{M}\) of \(\hat{\mathbf{V}} = \mathbf{V} \otimes_{\mathbf{K}} \hat{\mathbf{K}}\) with respect to \(\hat{\mathbf{A}}\) , is bijective and its inverse maps every lattice M' of \(\hat{\mathbf{V}}\) with respect to \(\hat{\mathbf{A}}\) to its intersection M' \(\cap\) V (V being canonically identified with a vector sub-K-space of \(\hat{\mathbf{V}}\) ).

If L is a free lattice of V, the lattices aL (for \(a \in A\) , \(a \neq 0\) ) form a fundamental system of neighbourhoods of 0 for a topology T on V (compatible with its A-module structure), which (when a basis of L over A is taken) is identified with the product topology on \(K^{n}\) ; by virtue of Proposition 2, a fundamental system of neighbourhoods of 0 for T also consists of all the lattices of V with respect to A; clearly \(\hat{V}\) is the completion of V with respect to T. Moreover, if m is the maximal ideal of A, the topology T induces on every lattice M of V with respect to A the m-adic topology since M is a finitely generated A-module (Chapter III, §3, no. 2, Theorem 2) and \(\hat{A}M\) is the completion of M with respect to this topology (Chapter III, §2, no. 12, Proposition 16); moreover, as M is open (and therefore closed) in V, \(\hat{A}M \cap V = M\) , which proves again the fact that \(\phi\) is injective (which follows directly from Proposition 4, (ii), since \(\hat{A}\) is a faithfully flat A-module). Finally, if \(M'\) is a lattice of \(\hat{V}\) with respect to \(\hat{A}\) , \(M = M' \cap V\) is a lattice of V with respect to A, for every element of \(\hat{A}\) is the product of an element of A and an invertible element of \(\hat{A}\) and hence it follows from Proposition

2 that there exist \(a, b\) in \(\mathbf{A} - \{0\}\) such that \(a\hat{\mathbf{A}}\mathbf{L} \subset \mathbf{M}' \subset b\hat{\mathbf{A}}\mathbf{L}\) , whence \(a\mathbf{L} \subset \mathbf{M}' \cap \mathbf{V} \subset b\mathbf{L}\) . Moreover \(\mathbf{M}'\) is open in \(\mathbf{V}\) and, as \(\mathbf{V}\) is dense in \(\hat{\mathbf{V}}\) , \(\mathbf{M}'\) is the completion of \(\mathbf{M}' \cap \mathbf{V} = \mathbf{M}\) ; this proves that \(\phi\) is surjective, whence the corollary.

Example (6) Let S be a multiplicative subset of A not containing 0; we apply Proposition 4 to \(B = S^{-1}A\) ; then L = K, \(BM = S^{-1}M\) ; hence \(S^{-1}M\) is a lattice of V with respect to \(S^{-1}A\) . Moreover:

PROPOSITION 5. Let V, W vector spaces of finite rank over K, M a lattice of V and N a lattice of W. If M is finitely generated, then (in the notation of Proposition 3):

\[
\mathbf {S} ^ {- 1} (\mathbf {N}: \mathbf {M}) = \mathbf {S} ^ {- 1} \mathbf {N}: \mathbf {S} ^ {- 1} \mathbf {M}\tag{1}
\]

in \(\operatorname{Hom}_{\mathbf{K}}(\mathbf{V},\mathbf{W})\)

Clearly the left hand side of (1) is contained in the right hand side. Conversely, let \(f \in S^{-1}N: S^{-1}M\) and let \((x_i)_{1 \leqslant i \leqslant n}\) be a system of generators of M. There exists \(s \in S\) such that \(f(x_i) \in s^{-1}N\) for all \(i\) and hence \(sf \in N: M\) , which proves the proposition.

